\origpage{85}{GRAPH MINORS. XIII}
\noindent automatically maximal, if we run the augmenting paths starting from $V(G_i)$. Thus this is easy.) We choose $\xi'$ minimum such that some $(A_i, B_i)$ has order $\xi'$, and we choose $i$ with $A_i$ maximal such that $(A_i, B_i)$ has order $\xi'$. We choose $v \in V(G_i)$. Then $v$ is the desired irrelevant vertex, by (6.1), since $(A_i, B_i)$ satisfies the conditions for $(A, B)$ in (6.1).

The algorithm involves computing $(A_i, B_i)$ for $\le p$ values of $i$; and each computation can be performed with a standard network flow algorithm in time $O(\varepsilon)$. Thus the algorithm has running time $O(\varepsilon)$. $\square$

\stmt
(6.2) will form one of the key steps of our algorithm for Folio, but there is also a peripheral application for it, to reduce our general problem to the case when $|E(G)| \le O(\nu)$ and, thus, to improve the overall running time of the main algorithm. In the remainder of this section we explain this.

If $G$ is a graph and $X \subseteq V(G)$, we denote by $N(X)$ the set of vertices in $V(G) - X$ with a neighbour in $X$. First, we need the following.

\stmt
(6.3) \textsc{Algorithm}.

\textit{Input:} An integer $k > 0$ and a non-null simple graph $G$ with $|E(G)| \ge k |V(G)|$.

\textit{Output:} A subset $X \subseteq V(G)$ such that

\hspace*{12pt}(i) the restriction of $G$ to $X$ is non-null and connected

\hspace*{12pt}(ii) $N(X)$ is non-null, and

\hspace*{12pt}(iii) every vertex of $N(X)$ has $\ge k$ neighbours in $N(X)$.

\textit{Running time:} $O(\nu\varepsilon)$.

\textit{Description:} (1) Choose a component $H$ of $G$ with $|E(H)| \ge k |V(H)|$, and choose $v_1 \in V(H)$.

\hspace*{12pt}(2) We perform a recursion. Set $X_1 = \{v_1\}$. At the start of the $i$ th step of the recursion we have a subset $X_i \subseteq V(H)$ with $|X_i| = i$ such that $X_i$ meets at most $|N(X_i)| + k(i - 1)$ edges, and such that the restriction of $G$ to $X_i$ is connected.

\hspace*{12pt}(3) If there exists $v \in N(X_i)$ with $< k$ neighbours in $N(X_i)$, we set $X_{i+1} = X_i \cup \{v\}$ and return for the next iteration. To check the inequality, let $v$ have $d_1$ neighbours in $N(X_i)$ and $d_2$ in $V(H) - (X_i \cup N(X_i))$. Then
\begin{equation*}
|N(X_{i+1})| = |N(X_i)| - 1 + d_2
\end{equation*}
and $X_{i+1}$ meets $d_1 + d_2 \le k - 1 + d_2$ edges not met by $X_i$. It follows that $X_{i+1}$ meets at most
\begin{equation*}
|N(X_i)| + k(i - 1) + k - 1 + d_2 = |N(X_{i+1})| + ki
\end{equation*}
edges in total, and so $X_{i+1}$ has the properties claimed.
